% =====================================================================
%  CIMENTA — Bitácora del proyecto
%  Fuente LaTeX. Compilar con:  tectonic BITACORA.tex
%  Para agregar una entrada nueva: copia el bloque \entrada{...} de
%  ejemplo al inicio de la sección de entradas (la más reciente arriba).
% =====================================================================
\documentclass[11pt]{article}

\usepackage{fontspec}            % XeLaTeX: UTF-8 y acentos nativos
\usepackage[spanish,es-noindentfirst]{babel}
\usepackage[a4paper,margin=2.4cm]{geometry}
\usepackage{xcolor}
\usepackage{titlesec}
\usepackage{enumitem}
\usepackage{booktabs}
\usepackage{fancyhdr}
\usepackage{parskip}
\usepackage{amssymb}             % \square para checkboxes
\usepackage[hidelinks]{hyperref}

% ---- Paleta de marca CIMENTA ----
\definecolor{azul}{HTML}{16324F}
\definecolor{acento}{HTML}{2D7DD2}
\definecolor{grafito}{HTML}{43484D}
\definecolor{niebla}{HTML}{AEB6BD}

\color{grafito}

% ---- Títulos ----
\titleformat{\section}{\Large\bfseries\color{azul}}{}{0pt}{}[\vspace{2pt}\color{niebla}\titlerule]
\titleformat{\subsection}{\large\bfseries\color{acento}}{}{0pt}{}
\titlespacing{\section}{0pt}{16pt}{6pt}
\titlespacing{\subsection}{0pt}{10pt}{2pt}

% ---- Encabezado / pie ----
\pagestyle{fancy}
\fancyhf{}
\renewcommand{\headrulewidth}{0.4pt}
\lhead{\small\color{azul}\textbf{CIMENTA}\,\textcolor{niebla}{·}\,Bitácora}
\rhead{\small\color{niebla}\thepage}
\renewcommand{\footrulewidth}{0pt}

% ---- Atajos ----
\newcommand{\codigo}[1]{\texttt{#1}}
\newcommand{\pend}{\item[\textcolor{acento}{$\square$}]}
\setlist[itemize]{leftmargin=1.4em,itemsep=2pt,topsep=2pt}

% =====================================================================
\begin{document}

% ---------- Portada / cabecera ----------
\begin{center}
  {\color{azul}\Huge\bfseries CIMENTA}\\[2pt]
  {\color{acento}\large Bitácora del proyecto}\\[6pt]
  {\color{niebla}\small Inteligencia inmobiliaria cuantitativa · Ciudad de México · 2026}
\end{center}
\vspace{6pt}
\hrule height 1pt
\vspace{10pt}

\noindent Registro cronológico de todo lo que hacemos en el proyecto.
La entrada más reciente va arriba.

\textbf{Convención de cada entrada:} fecha y responsable · qué se hizo ·
decisiones / notas · pendientes / siguiente paso.

% =====================================================================
%  ENTRADAS (la más reciente arriba)
% =====================================================================

\section*{2026-06-12 — Arranque del proyecto (Etapa 0: setup)}
\addcontentsline{toc}{section}{2026-06-12 — Arranque del proyecto}

\subsection*{Qué se hizo}
\begin{itemize}
  \item Estructura base del repositorio: \codigo{app/}, \codigo{data/},
        \codigo{notebooks/}, \codigo{test/}, \codigo{src/}, \codigo{docs/},
        más \codigo{README.md}, \codigo{requirements.txt}, \codigo{.gitignore}
        y esta bitácora.
  \item Carpetas vacías versionadas con \codigo{.gitkeep}.
  \item Entorno de Anaconda \textbf{\codigo{cimenta\_env}} con \textbf{Python 3.12}
        (vacío, sin dependencias todavía).
  \item Repositorio Git con tres ramas: \codigo{master}, \codigo{dev},
        \codigo{feature/dev}.
  \item Remoto configurado a
        \codigo{https://github.com/danielomar14/cimenta\_repo.git} y las tres
        ramas subidas.
\end{itemize}

\subsection*{Decisiones / notas}
\begin{itemize}
  \item Flujo de trabajo: desarrollo en \codigo{feature/dev}
        $\rightarrow$ pruebas en \codigo{dev}
        $\rightarrow$ estable en \codigo{master}.
  \item El entorno se mantiene vacío a propósito; las dependencias se agregarán
        etapa por etapa.
  \item Datos (\codigo{data/}) excluidos del control de versiones vía
        \codigo{.gitignore}.
  \item La bitácora se lleva en LaTeX (\codigo{BITACORA.tex}) y se renderiza a
        PDF (\codigo{BITACORA.pdf}) con \codigo{tectonic}.
\end{itemize}

\subsection*{Bloqueo detectado — Catálogo de Etapas ilegible}
\begin{itemize}
  \item El archivo \codigo{docs/context/CIMENTA-Catalogo-de-Etapas.pdf}
        (236 págs.) está \textbf{corrupto}: todos los bytes $\geq$ \codigo{0x80}
        de sus streams comprimidos fueron sustituidos por el carácter de
        reemplazo Unicode U+FFFD (\codigo{ef bf bd}) — 823{,}069 ocurrencias.
        Es resultado de un round-trip UTF-8 que destruyó la data binaria.
        \textbf{No es recuperable}: ni texto ni render.
  \item Los demás PDFs de \codigo{docs/context/} están sanos: Plan-Maestro-v2
        (97p), Plan-Maestro (43p), Bitácora-del-Jurado (120p), Propuesta-Inicial
        (12p), Riesgos-y-Debilidades (17p).
  \item \textbf{Acción requerida:} reexportar / volver a subir una copia limpia
        del Catálogo de Etapas para poder avanzar etapa por etapa.
\end{itemize}

\subsection*{Pendientes / siguiente paso}
\begin{itemize}
  \pend Obtener copia limpia de \codigo{CIMENTA-Catalogo-de-Etapas.pdf}.
  \pend Leer el catálogo y definir la Etapa 1.
\end{itemize}

\end{document}
